\documentclass[11pt]{article}

\usepackage[utf8]{inputenc}
\usepackage{graphicx}
\usepackage[dvipsnames]{xcolor}
\usepackage{color}
\usepackage[margin=1in]{geometry}
\usepackage{hyperref}
\usepackage{tikz}
\usepackage{amsmath}
\usepackage{commath}
\usepackage{relsize}
\usepackage{centernot}
\usepackage{amssymb}
\usepackage{amsfonts}
\usepackage{graphicx}
\usepackage{textcomp}
\usepackage{gensymb}
\usepackage{enumitem}
\usepackage{siunitx}
\usepackage{fancyhdr}
\usepackage{floatrow}
\usepackage{wrapfig}
\usepackage[labelfont=bf]{caption}
\usepackage{mathtools}
\usepackage{pgfplots}
\usepackage{multicol}
\usepackage[misc]{ifsym}
\usepackage{tabularx}
\usepackage{empheq}
\usepackage{tkz-tab}
\usepackage{xpatch}
\usepackage{caption}
\usepackage{subcaption}
\usepackage{tabularray}
\usepackage{esint}
\usepackage{amsthm}
\usepackage{thmtools}
\RequirePackage[most]{tcolorbox}

\swapnumbers

\makeatletter
\declaretheoremstyle[
  headfont= \sffamily\bfseries\color{cyan!50!black},
  headpunct={\sffamily\bfseries.},
  postheadspace=0.5em,
  notefont=\sffamily\bfseries,
  headformat=\NAME\NUMBER\let\thmt@space\@empty\NOTE,
  bodyfont=\mdseries\slshape,
  spaceabove=12pt,spacebelow=12pt,
]{thmstyle}
\makeatother

\theoremstyle{thmstyle}
\declaretheorem[name=Definición]{theorem}
\newtheorem{definition}[theorem]{Definición}

\definecolor{lightcyan}{rgb}{0.88, 1.0, 1.0}
\colorlet{mythmback}{lightcyan!40!white}

\tcolorboxenvironment{theorem}{
  colback=mythmback,coltitle=blue,colframe=mythmback,
  left=0pt,right=0pt,
  top=0pt,bottom=0pt,
  before skip=10pt, after skip=10pt,
}

\xpatchcmd{\tkzTabLine}{$0$}{$\bullet$}{}{}
\tikzset{t style/.style={style=solid}}
\pgfplotsset{compat=newest}

\newtheorem{manualtheoreminner}{}
\newenvironment{manualtheorem}[1]{%
  \renewcommand\themanualtheoreminner{#1}%
  \manualtheoreminner
}{\endmanualtheoreminner}

\renewcommand{\pd}[2]{\frac{\partial #1}{ \partial #2}}
\newcommand{\td}[2]{\frac{\mathrm{d} #1}{ \mathrm{d} #2}}
\newcommand{\pdd}[2]{\frac{\partial^2 #1}{ \partial #2 ^2}}
\newcommand{\pddc}[3]{\frac{\partial^2 #1}{ \partial #2 \partial #3}}
\newcommand{\solution}[0]{\large{\textbf{Solución}}\normalsize}

\def\sca   #1{\mbox{\rm{#1}}{}}
\def\mat   #1{\mbox{\boldmath $\mathsf #1$}}
\def\vec   #1{\mbox{\boldmath $#1$}{}}
\def\ten   #1{\mbox{\boldmath $#1$}{}}
\def\ltr   #1{\mbox{\sf{#1}}}
\def\bltr  #1{\mbox{\sffamily{\bfseries{{#1}}}}}

\DeclareMathOperator{\dive}{div}
\DeclareMathOperator{\trace}{trace}
\DeclareMathOperator{\tr}{tr}
\DeclareMathOperator{\symm}{symm}
\DeclareMathOperator{\sk}{skew}
\DeclareMathOperator{\grad}{grad}
\DeclareMathOperator{\Grad}{Grad}
\DeclareMathOperator{\curl}{curl}
\DeclareMathOperator{\Curl}{Curl}
\DeclareMathOperator*{\argmin}{\arg\!\min}

\def\dx{\mbox{\(\,\mathrm{d}x\)}}
\def\ds{\mbox{\(\,\mathrm{d}s\)}}
\def\dS{\mbox{\(\,\mathrm{d}S\)}}
\def\dV{\mbox{\(\,\mathrm{d}V\)}}
\def\dv{\mbox{\(\,\mathrm{d}v\)}}
\def\R{\mathbb R}
\def\C{\mathbb C}
\def\N{\mathbb N}
\def\Q{\mathbb Q}
\def\Z{\mathbb Z}

\usepackage{mdframed}
\mdfdefinestyle{MyFrame}{%
    linecolor=black,
    linewidth=1pt,
    outerlinewidth=1pt,
    roundcorner=20pt,
    innertopmargin=\baselineskip,
    innerbottommargin=\baselineskip,
    innerrightmargin=20pt,
    innerleftmargin=20pt,
    splittopskip=2\baselineskip,
    backgroundcolor=gray!50!white}

\setlength\textwidth{7in}
\setlength\parindent{0pt}
\oddsidemargin=-0.5cm
\pagestyle{fancy}
\newcommand*{\vertbar}{\rule[-1ex]{0.5pt}{2.5ex}}
\setlength{\headheight}{13.59999pt}

\renewcommand{\headrulewidth}{0.1pt}
\renewcommand{\footrulewidth}{0.1pt}
\renewcommand{\figurename}{Figura}
\renewcommand\refname{Referencias}
\renewcommand{\thesection}{\arabic{section}.}
\renewcommand{\thesubsection}{\thesection\arabic{subsection}.}
\renewcommand{\labelitemi}{{\raisebox{.3\height}{\scalebox{0.6}{$\blacklozenge$}}}}
\renewcommand*\contentsname{\'Indice}

\newcommand{\p}{\vspace{10pt}}

\AtBeginDocument{\vspace*{-3.2\baselineskip}}
\textheight=665pt

\lhead{\scshape Pontificia Universidad Cat\'olica de Chile}
\rhead{\scshape IMC UC}
\cfoot{\thepage}

\usepackage{footnote}
\usepackage{etoolbox}
\newtoggle{indefintion}
\togglefalse{indefintion}
\pretocmd{\footnote}{\iftoggle{indefintion}{\stepcounter{footnote}}{\relax}}{}{}

\begin{document}
    \begin{flushleft}
        \small
        \vspace{1pt}
        \textbf{IMT3130} \hfill
        \textbf{27/03/2026}
    \end{flushleft}
    \begin{center}
        \LARGE\textbf{Ayudant\'ia 3}\\
        \vspace{3pt}
        \normalsize\textbf{Diferencias finitas en 2D, stepping temporal y estabilidad}\\
        \normalsize Ayudante: Basti\'an Herrera \Letter{\hspace{1pt}: \texttt{biherrera@uc.cl}}
        \rule{\textwidth}{0.05mm}
    \end{center}

\section{Problemas}

\begin{enumerate}[label=\textbf{\arabic*.}, leftmargin=*]

    \item Considere el problema de Poisson
    \[
    \begin{aligned}
        -\Delta u &= f && \text{en } \Omega=(0,1)^2, \\
        u &= 0 && \text{en } \partial\Omega.
    \end{aligned}
    \]
    Sobre la malla uniforme $x_i=ih$, $y_j=jh$, con $h=1/N$, sea $u_{ij}\approx u(x_i,y_j)$ en nodos interiores.
    \begin{enumerate}[label=(\alph*)]
        \item Escriba el esquema de 5 puntos y el stencil asociado.
        \item Usando la soluci\'on manufacturada $u(x,y)=\sin(\pi x)\sin(\pi y)$, determine $f$ y el orden esperado del m\'etodo.
        \item Implemente el esquema para $h=1/8,1/16,1/32,1/64$, reporte el error máximo en la malla $E_h=\max_{i,j}\abs{u(x_i,y_j)-u_{ij}}$, y estime la tasa emp\'irica de convergencia.
    \end{enumerate}

    % \item Considere la ecuaci\'on del calor unidimensional
    % \[
    % \pd{u}{t}-\pdd{u}{x}=0, \qquad (x,t)\in(0,1)\times(0,T],
    % \]
    % con condiciones de borde homog\'eneas y condici\'on inicial $u_0$. Sea $x_j=jh$, $t^n=n\Delta t$ y $U_j^n\approx u(x_j,t^n)$. Considere el esquema $\theta$
    % \[
    % \frac{U_j^{n+1}-U_j^n}{\Delta t}
    % =
    % \theta \frac{U_{j+1}^{n+1}-2U_j^{n+1}+U_{j-1}^{n+1}}{h^2}
    % +(1-\theta)\frac{U_{j+1}^{n}-2U_j^{n}+U_{j-1}^{n}}{h^2}.
    % \]
    % \begin{enumerate}[label=(\alph*)]
    %     \item Identifique los casos $\theta=0$, $\theta=1$ y $\theta=1/2$.
    %     \item Muestre que el error de truncamiento local satisface
    %     \[
    %     \tau_j^n=O(h^2+\Delta t)
    %     \quad \text{si } \theta\neq \frac12,
    %     \qquad
    %     \tau_j^n=O(h^2+\Delta t^2)
    %     \quad \text{si } \theta=\frac12.
    %     \]
    %     \item Según estos resultados, explique la ventaja pr\'actica de Crank-Nicolson frente a Euler expl\'icito e impl\'icito.
    % \end{enumerate}

    \item Considere la ecuaci\'on de advecci\'on peri\'odica
    \[
    \pd{u}{t}+\pd{u}{x}=0, \qquad x\in[0,1], \quad t>0,
    \]
    con condici\'on peri\'odica y un dato inicial suave. Estudiamos el esquema
    \[
    \frac{U_j^{n+1}-U_j^n}{\Delta t} + \frac{U_{j+1}^{n}-U_{j-1}^{n}}{2h}=0.
    \]
    \begin{enumerate}[label=(\alph*)]
        \item Determine su orden de consistencia.
        \item Realice un an\'alisis de von Neumann y obtenga el factor de amplificaci\'on $g(\xi)$.
        \item Pruebe que el esquema es inestable para todo $\lambda=\Delta t/h>0$ y explique por qu\'e este ejemplo muestra que consistencia no implica convergencia.
    \end{enumerate}

    \item Considere nuevamente la ecuaci\'on de advecci\'on peri\'odica
    \[
    \pd{u}{t}+a\pd{u}{x}=0, \qquad a\in\R..
    \]
    \begin{enumerate}[label=(\alph*)]
        \item Sea $\nu=a\Delta t/h$. Derive el esquema de Lax-Wendroff
        \[
        U_j^{n+1}
        =
        U_j^n
        -\frac{a\Delta t}{2h}(U_{j+1}^n-U_{j-1}^n)
        +\frac{a^2\Delta t^2}{2h^2}(U_{j+1}^n-2U_j^n+U_{j-1}^n)
        \]
        usando la ecuación y las expansiones de Taylor en tiempo y espacio. 
        \item Muestre que el error local de truncamiento es de orden $2$ y que su t\'ermino principal est\'a dado por una tercera derivada espacial, es decir,
        \[
        \tau_j^n = \frac{a h^2}{6}(1-\nu^2)\,u_{xxx}(x_j,t^n) + O(h^3),
        \]
        asumiendo $\nu$ fijo.
        \item Realice un an\'alisis de von Neumann y pruebe que el esquema es estable si y solo si $\abs{\nu}\le 1$.
        \item Usando el teorema de equivalencia de Lax, explique qu\'e se concluye sobre la convergencia del esquema cuando $\abs{\nu}\le 1$.
    \end{enumerate}

\end{enumerate}

\end{document}
